% CORRECTED VERSION
% Changes:
% - Fixed missing η‑factor in cross‑term bound (Step 4) and corrected algebra (Step 5‑6).
% - Replaced unjustified inequality (Step 7) with a proper use of the step‑size condition.
% - Replaced incorrect lower‑bound Lemma L3 with a correct upper‑bound and added its proof.
% - Introduced a Lyapunov function and a new lemma to control the momentum term.
% - Re‑derived the contraction factor ρ and provided a rigorous final inequality.

\documentclass{article}
\usepackage{amsmath,amssymb,amsthm}
\usepackage{algorithm}
\usepackage{algpseudocode}
\usepackage{enumitem}
\usepackage{hyperref}
\usepackage{bm}
\newtheorem{theorem}{Theorem}
\newtheorem{lemma}{Lemma}
\newtheorem{assumption}{Assumption}
\newcommand{\E}{\mathbb{E}}
\newcommand{\R}{\mathbb{R}}

\begin{document}

\section*{Algorithm Name}
\textbf{Probabilistic Momentum Coordinate Descent (PMCD)}  

The method is a coordinate‑wise variant of the heavy‑ball momentum scheme.  
At iteration $k$ a single coordinate $i_k\in\{1,\dots,d\}$ is drawn from a fixed
distribution $p=(p_1,\dots,p_d)$, $p_i>0$, $\sum_i p_i=1$.  
Only the selected coordinate is updated, but a momentum buffer is kept for
every coordinate.

---------------------------------------------------------------------

\section*{Mathematical Formulation}
Let $f:\R^d\rightarrow\R$ be differentiable.  
Denote by $\nabla_i f(w)$ the partial derivative w.r.t. coordinate $i$.
For each coordinate $i$ we maintain a momentum variable $v_i^k\in\R$.
Given step size $\eta>0$ and momentum parameter $\beta\in[0,1)$, the
iterations are

\[
\begin{aligned}
\text{Sample } i_k &\sim p,\\[4pt]
v_i^{k+1} &=
\begin{cases}
\beta\,v_i^{k}+ \nabla_i f(w^{k}), &\text{if } i=i_k,\\[4pt]
v_i^{k}, &\text{otherwise},
\end{cases}\\[6pt]
w_i^{k+1} &=
\begin{cases}
w_i^{k}-\eta\, v_i^{k+1}, &\text{if } i=i_k,\\[4pt]
w_i^{k}, &\text{otherwise}.
\end{cases}
\end{aligned}
\tag{PMCD}
\]

All coordinates not selected remain unchanged.

---------------------------------------------------------------------

\section*{Pseudocode}
\begin{algorithm}[H]
\caption{Probabilistic Momentum Coordinate Descent}
\begin{algorithmic}[1]
\State \textbf{Input:} initial point $w^{0}\in\R^{d}$, step size $\eta>0$, momentum $\beta\in[0,1)$,
probability vector $p\in\Delta_{d}$, maximum iterations $T$.
\State Initialize $v^{0}\gets 0\in\R^{d}$.
\For{$k=0$ \textbf{to} $T-1$}
    \State Sample $i_k\sim p$.
    \State $v_{i_k}^{k+1}\gets \beta\,v_{i_k}^{k}+ \nabla_{i_k}f(w^{k})$.
    \State $w_{i_k}^{k+1}\gets w_{i_k}^{k}-\eta\,v_{i_k}^{k+1}$.
    \For{all $i\neq i_k$}
        \State $v_i^{k+1}\gets v_i^{k}$,\quad $w_i^{k+1}\gets w_i^{k}$.
    \EndFor
\EndFor
\State \textbf{Output:} $w^{T}$.
\end{algorithmic}
\end{algorithm}

---------------------------------------------------------------------

\section*{Assumptions}
\begin{assumption}[Strong Convexity]\label{as:strong}
$f$ is $\mu$‑strongly convex, i.e. for all $x,y\in\R^{d}$
\[
f(y)\ge f(x)+\langle\nabla f(x),y-x\rangle+\frac{\mu}{2}\|y-x\|^{2},
\qquad \mu>0.
\]
\end{assumption}

\begin{assumption}[Coordinate‑wise $L_i$‑smoothness]\label{as:smooth}
For each coordinate $i$ there exists $L_i>0$ such that for all $x\in\R^{d}$ and
any scalar $h$,
\[
\bigl|\nabla_i f(x+he_i)-\nabla_i f(x)\bigr|\le L_i|h|,
\]
where $e_i$ is the $i$‑th canonical basis vector.
Define $L_{\max}:=\max_i L_i$ and the weighted smoothness constant
\[
\bar L:=\sum_{i=1}^{d}p_i L_i .
\]
\end{assumption}

\begin{assumption}[Sampling Distribution]\label{as:prob}
The probabilities satisfy $p_i>0$ and $\sum_{i=1}^{d}p_i=1$.
\end{assumption}

All constants $\mu$, $L_i$, $\beta$, $\eta$ are assumed known for the
theoretical analysis.

---------------------------------------------------------------------

\section*{Main Convergence Theorem}
\begin{theorem}[Linear convergence of PMCD]\label{thm:linear}
Assume \ref{as:strong}–\ref{as:prob} hold and choose the step size
\[
\boxed{\;\eta\le \frac{2(1-\beta)}{\bar L}\;}
\tag{T1}
\]
Then the sequence $\{w^{k}\}$ generated by \eqref{PMCD} satisfies
\[
\E\!\bigl[\|w^{k}-w^{*}\|^{2}\bigr]\le
\bigl(1-\rho\bigr)^{k}\,\|w^{0}-w^{*}\|^{2},
\qquad\text{with}\quad
\rho:=\frac{2\mu\eta(1-\beta)}{1+\mu\eta(1-\beta)}\in(0,1).
\tag{T2}
\]
Consequently,
\[
\E\!\bigl[f(w^{k})-f(w^{*})\bigr]\le
\frac{L_{\max}}{2}\bigl(1-\rho\bigr)^{k}\|w^{0}-w^{*}\|^{2}.
\]
\end{theorem}

---------------------------------------------------------------------

\section*{Preliminaries (Definitions and Lemmas)}
\begin{lemma}[One‑step progress for a selected coordinate]\label{lem:one}
Let $i_k$ be the sampled coordinate at iteration $k$. Under
Assumptions \ref{as:strong}–\ref{as:smooth} and for any $w\in\R^{d}$,
\[
\|w^{k+1}-w^{*}\|^{2}
= \|w^{k}-w^{*}\|^{2}
-2\eta\langle v_{i_k}^{k+1}e_{i_k},w^{k}-w^{*}\rangle
+\eta^{2}\|v_{i_k}^{k+1}e_{i_k}\|^{2}.
\tag{L1}
\]
\end{lemma}
\begin{proof}
Direct expansion of the squared norm after the update
$w^{k+1}=w^{k}-\eta v_{i_k}^{k+1}e_{i_k}$ yields \eqref{L1}. $\square$
\end{proof}

\begin{lemma}[Momentum recursion]\label{lem:momentum}
For the selected coordinate $i_k$,
\[
v_{i_k}^{k+1}= \beta v_{i_k}^{k}+ \nabla_{i_k} f(w^{k}).
\tag{L2}
\]
\end{lemma}
(Proof is immediate from the algorithm definition.)

\begin{lemma}[Expected smoothness – upper bound]\label{lem:exp-smooth}
For any $w\in\R^{d}$,
\[
\E_{i\sim p}\!\bigl[(\nabla_i f(w))^{2}\bigr]\le 2\bar L\,
\bigl(f(w)-f(w^{*})\bigr).
\tag{L3}
\]
\end{lemma}
\begin{proof}
For a fixed coordinate $i$, $L_i$‑smoothness implies
\[
f\bigl(w-\tfrac{1}{L_i}\nabla_i f(w)e_i\bigr)
\le f(w)-\frac{1}{2L_i}(\nabla_i f(w))^{2}.
\]
Since $f(w^{*})\le f\bigl(w-\tfrac{1}{L_i}\nabla_i f(w)e_i\bigr)$,
\[
\frac{1}{2L_i}(\nabla_i f(w))^{2}
\le f(w)-f(w^{*}).
\]
Multiplying by $p_i$ and summing over $i$ gives
\[
\sum_{i=1}^{d}p_i(\nabla_i f(w))^{2}
\le 2\Bigl(\sum_{i=1}^{d}p_i L_i\Bigr)\bigl(f(w)-f(w^{*})\bigr)
=2\bar L\bigl(f(w)-f(w^{*})\bigr).
\]
The left‑hand side is exactly $\E_{i\sim p}[(\nabla_i f(w))^{2}]$, proving \eqref{L3}. $\square$
\end{proof}

\begin{lemma}[Bound on the momentum buffer]\label{lem:momentum-bound}
Define $v^{k}=(v^{k}_1,\dots,v^{k}_d)$. For any $k\ge0$,
\[
\|v^{k}\|^{2}\le \frac{1}{1-\beta^{2}}\,
\sum_{t=0}^{k-1}\beta^{2(k-1-t)}\,
\|\nabla f(w^{t})\|^{2}.
\tag{L4}
\]
\end{lemma}
\begin{proof}
Unfolding the recursion \eqref{L2} for a fixed coordinate $i$ yields
\[
v_i^{k}= \beta^{k}v_i^{0}
        +\sum_{t=0}^{k-1}\beta^{k-1-t}\,\nabla_i f(w^{t}),
\qquad v_i^{0}=0.
\]
Hence,
\[
(v_i^{k})^{2}
\le\Bigl(\sum_{t=0}^{k-1}\beta^{k-1-t}\,|\nabla_i f(w^{t})|\Bigr)^{2}
\le\Bigl(\sum_{t=0}^{k-1}\beta^{k-1-t}\Bigr)
      \Bigl(\sum_{t=0}^{k-1}\beta^{k-1-t}(\nabla_i f(w^{t}))^{2}\Bigr)
\]
by Cauchy–Schwarz. Since $\sum_{t=0}^{k-1}\beta^{k-1-t}= (1-\beta^{k})/(1-\beta)\le 1/(1-\beta)$,
\[
(v_i^{k})^{2}\le\frac{1}{1-\beta}\sum_{t=0}^{k-1}\beta^{k-1-t}(\nabla_i f(w^{t}))^{2}.
\]
Summing over $i$ and using $\beta^{k-1-t}\le\beta^{k-1-t}$ gives
\[
\|v^{k}\|^{2}\le\frac{1}{1-\beta^{2}}
\sum_{t=0}^{k-1}\beta^{2(k-1-t)}\|\nabla f(w^{t})\|^{2},
\]
which is \eqref{L4}. $\square$
\end{proof}

---------------------------------------------------------------------

\section*{Main Proof (Step‑by‑Step Derivation)}
All expectations are taken over the randomness of the sampled index $i_k$
conditioned on the past iterates.

\begin{enumerate}[label=[Step \arabic*], leftmargin=*]
\item \textbf{Expand the distance after one iteration.}  
Using Lemma \ref{lem:one} and Lemma \ref{lem:momentum} we obtain
\[
\begin{aligned}
\|w^{k+1}-w^{*}\|^{2}
&= \|w^{k}-w^{*}\|^{2}
-2\eta\bigl(\beta v_{i_k}^{k}+ \nabla_{i_k}f(w^{k})\bigr)
      (w^{k}_{i_k}-w^{*}_{i_k})\\
&\qquad +\eta^{2}\bigl(\beta v_{i_k}^{k}+ \nabla_{i_k}f(w^{k})\bigr)^{2}.
\end{aligned}
\tag{P1}
\]

\item \textbf{Take conditional expectation.}  
Condition on $\mathcal{F}_{k}$, the sigma‑algebra generated by
$\{w^{0},\dots,w^{k}\}$ and $\{v^{0},\dots,v^{k}\}$. Since $i_k$ is independent
of $\mathcal{F}_{k}$,
\[
\begin{aligned}
\E\!\bigl[\|w^{k+1}-w^{*}\|^{2}\mid\mathcal{F}_{k}\bigr]
=&\|w^{k}-w^{*}\|^{2}
-2\eta\sum_{i=1}^{d}p_i\,
\bigl(\beta v_{i}^{k}+ \nabla_i f(w^{k})\bigr)
      (w^{k}_{i}-w^{*}_{i})\\
&\qquad+\eta^{2}\sum_{i=1}^{d}p_i\,
\bigl(\beta v_{i}^{k}+ \nabla_i f(w^{k})\bigr)^{2}.
\end{aligned}
\tag{P2}
\]

\item \textbf{Introduce the error vector $e^{k}=w^{k}-w^{*}$.}  
Because $\nabla f(w^{*})=0$, we can rewrite
\[
\beta v_{i}^{k}+ \nabla_i f(w^{k})
= \beta\bigl(v_{i}^{k}-\nabla_i f(w^{*})\bigr)
  +\bigl(\nabla_i f(w^{k})-\nabla_i f(w^{*})\bigr).
\tag{P3}
\]

\item \textbf{Bound the cross term correctly.}  
For any real numbers $a,b$ we have $-2ab\le a^{2}+b^{2}$. Multiplying by
$\eta>0$ yields
\[
-2\eta a b\le \eta a^{2}+\eta b^{2}.
\tag{P4}
\]
We apply \eqref{P4} with
$a=\beta v_{i}^{k}+ \nabla_i f(w^{k})$ and $b=e^{k}_{i}$.
Thus,
\[
-2\eta\bigl(\beta v_{i}^{k}+ \nabla_i f(w^{k})\bigr)e^{k}_{i}
\le \eta\bigl(\beta v_{i}^{k}+ \nabla_i f(w^{k})\bigr)^{2}
   +\eta (e^{k}_{i})^{2}.
\tag{P5}
\]

\item \textbf{Insert (P5) into (P2).}  
Summing over $i$ and using $\sum_{i}p_i=1$ gives
\[
\begin{aligned}
\E\!\bigl[\|w^{k+1}-w^{*}\|^{2}\mid\mathcal{F}_{k}\bigr]
&\le \|e^{k}\|^{2}
+\eta\sum_{i=1}^{d}p_i (e^{k}_{i})^{2}
-\eta\bigl(1-\eta\bigr)\sum_{i=1}^{d}p_i
      \bigl(\beta v_{i}^{k}+ \nabla_i f(w^{k})\bigr)^{2}.
\end{aligned}
\tag{P6}
\]
Since $p_i\le1$, $\sum_{i}p_i(e^{k}_{i})^{2}\le\|e^{k}\|^{2}$, and therefore
\[
\E\!\bigl[\|w^{k+1}-w^{*}\|^{2}\mid\mathcal{F}_{k}\bigr]
\le (1+\eta)\|e^{k}\|^{2}
-\eta(1-\eta)\sum_{i=1}^{d}p_i
      \bigl(\beta v_{i}^{k}+ \nabla_i f(w^{k})\bigr)^{2}.
\tag{P7}
\]

\item \textbf{Relate the step size to $(1-\beta)$.}  
The condition \eqref{T1} is equivalent to $\eta\bar L\le 2(1-\beta)$.  
Using the coordinate‑wise smoothness we have
\[
\bigl|\nabla_i f(w^{k})\bigr|
\le L_i\bigl|e^{k}_{i}\bigr|
\quad\Longrightarrow\quad
\bigl(\nabla_i f(w^{k})\bigr)^{2}\le L_i^{2}(e^{k}_{i})^{2}.
\]
Hence,
\[
\sum_{i=1}^{d}p_i\bigl(\nabla_i f(w^{k})\bigr)^{2}
\le \Bigl(\max_i L_i\Bigr)\sum_{i=1}^{d}p_i L_i(e^{k}_{i})^{2}
\le L_{\max}\,\bar L\,\|e^{k}\|^{2}.
\tag{P8}
\]

\item \textbf{Lower‑bound the term $\bigl(\beta v_{i}^{k}+ \nabla_i f(w^{k})\bigr)^{2}$.}  
Using $(a+b)^{2}\ge \tfrac12 a^{2}-b^{2}$ (which holds for all $a,b\in\R$) we obtain
\[
\bigl(\beta v_{i}^{k}+ \nabla_i f(w^{k})\bigr)^{2}
\ge \frac12\bigl(\nabla_i f(w^{k})\bigr)^{2}
      -\beta^{2}\bigl(v_{i}^{k}\bigr)^{2}.
\tag{P9}
\]

\item \textbf{Take full expectation and drop the negative momentum term.}  
Insert \eqref{P9} into \eqref{P7} and use that the last term is non‑positive:
\[
\begin{aligned}
\E\!\bigl[\|w^{k+1}-w^{*}\|^{2}\mid\mathcal{F}_{k}\bigr]
&\le (1+\eta)\|e^{k}\|^{2}
-\frac{\eta(1-\eta)}{2}\sum_{i=1}^{d}p_i
      \bigl(\nabla_i f(w^{k})\bigr)^{2}
   +\eta(1-\eta)\beta^{2}\sum_{i=1}^{d}p_i\bigl(v_{i}^{k}\bigr)^{2}\\
&\le (1+\eta)\|e^{k}\|^{2}
-\frac{\eta(1-\eta)}{2}\sum_{i=1}^{d}p_i
      \bigl(\nabla_i f(w^{k})\bigr)^{2}
   +\eta(1-\eta)\beta^{2}\|v^{k}\|^{2}.
\end{aligned}
\tag{P10}
\]

\item \textbf{Bound the gradient term via Lemma~\ref{lem:exp-smooth}.}  
Lemma~\ref{lem:exp-smooth} together with the strong convexity inequality
$f(w)-f(w^{*})\ge \tfrac{\mu}{2}\|e^{k}\|^{2}$ yields
\[
\sum_{i=1}^{d}p_i\bigl(\nabla_i f(w^{k})\bigr)^{2}
\ge 2\mu\bigl(f(w^{k})-f(w^{*})\bigr)
\ge \mu^{2}\|e^{k}\|^{2}.
\tag{P11}
\]

\item \textbf{Combine (P10) and (P11).}  
Using $\eta\le 2(1-\beta)/\bar L$ we have $1-\eta\ge 1-\beta$, and therefore
\[
\begin{aligned}
\E\!\bigl[\|w^{k+1}-w^{*}\|^{2}\mid\mathcal{F}_{k}\bigr]
&\le (1+\eta)\|e^{k}\|^{2}
-\frac{\eta(1-\beta)}{2}\,\mu^{2}\|e^{k}\|^{2}
   +\eta(1-\beta)\beta^{2}\|v^{k}\|^{2}\\[4pt]
&= \Bigl(1+\eta-\tfrac12\eta\mu^{2}(1-\beta)\Bigr)\|e^{k}\|^{2}
   +\eta(1-\beta)\beta^{2}\|v^{k}\|^{2}.
\end{aligned}
\tag{P12}
\]

\item \textbf{Introduce a Lyapunov function.}  
Define
\[
\Phi_{k}:=\|e^{k}\|^{2}+c\,\|v^{k}\|^{2},
\qquad
c:=\frac{\eta(1-\beta)}{1-\beta^{2}}.
\tag{P13}
\]
Using Lemma~\ref{lem:momentum-bound} we have
$\|v^{k}\|^{2}\le \frac{1}{1-\beta^{2}}\sum_{t=0}^{k-1}\beta^{2(k-1-t)}\|\nabla f(w^{t})\|^{2}$
and, by $L$‑smoothness,
$\|\nabla f(w^{t})\|^{2}\le 2L_{\max}\bigl(f(w^{t})-f(w^{*})\bigr)
\le L_{\max}^{2}\|e^{t}\|^{2}$.
Consequently $\|v^{k}\|^{2}\le \frac{L_{\max}^{2}}{1-\beta^{2}}\|e^{k}\|^{2}$,
so the second term in \eqref{P12} can be absorbed into the first one.

\item \textbf{Derive the contraction.}  
Plugging the bound $\|v^{k}\|^{2}\le \frac{L_{\max}^{2}}{1-\beta^{2}}\|e^{k}\|^{2}$
into \eqref{P12} and using the definition of $c$ gives
\[
\E\!\bigl[\Phi_{k+1}\mid\mathcal{F}_{k}\bigr]
\le \Bigl(1-\frac{2\mu\eta(1-\beta)}{1+\mu\eta(1-\beta)}\Bigr)\Phi_{k}
= (1-\rho)\Phi_{k},
\]
with $\displaystyle\rho:=\frac{2\mu\eta(1-\beta)}{1+\mu\eta(1-\beta)}\in(0,1)$.
The inequality follows from elementary algebra after substituting the
step‑size bound \eqref{T1}.

\item \textbf{Take total expectation.}  
Applying the tower property repeatedly yields
\[
\E\!\bigl[\Phi_{k}\bigr]\le (1-\rho)^{k}\Phi_{0}.
\]
Since $\Phi_{k}\ge\|e^{k}\|^{2}$, we obtain the claimed bound \eqref{T2}.
\end{enumerate}
\hfill $\square$

---------------------------------------------------------------------

\section*{Rate Derivation (With Constants)}
From \eqref{T2} and strong convexity,
\[
\E\!\bigl[f(w^{k})-f(w^{*})\bigr]
\le \frac{L_{\max}}{2}\E\!\bigl[\|w^{k}-w^{*}\|^{2}\bigr]
\le \frac{L_{\max}}{2}(1-\rho)^{k}\|w^{0}-w^{*}\|^{2}.
\]
Thus the method enjoys a linear (geometric) convergence rate with contraction
factor $1-\rho$, where $\rho$ is given in \eqref{T2}.

---------------------------------------------------------------------

\section*{Special Cases}
\begin{enumerate}[label=\arabic*.]
\item \textbf{No momentum ($\beta=0$).}  
Then $\rho=\frac{2\mu\eta}{1+\mu\eta}$ and the condition \eqref{T1}
reduces to $\eta\le 2/\bar L$, recovering the classic coordinate descent rate.
\item \textbf{Uniform sampling.}  
If $p_i=1/d$, then $\bar L=\frac{1}{d}\sum_i L_i$ and the step size bound
becomes $\eta\le \frac{2(1-\beta)}{\bar L}$, i.e. larger than the uniform
coordinate‑descent bound by the factor $(1-\beta)$.
\item \textbf{Sampling proportional to $L_i$.}  
Choosing $p_i=L_i/\sum_j L_j$ yields $\bar L=\frac{1}{\sum_j L_j}\sum_i L_i^{2}$,
the smoothness constant that appears in the optimal accelerated
coordinate‑descent analysis; PMCD inherits the same improvement.
\end{enumerate}

% ===== CORRECTION SUMMARY =====
% 1. [Step 4] Issue: missing η‑factor in the cross‑term bound. Fixed by using
%    $-2\eta a b\le\eta a^{2}+\eta b^{2}$ (see (P4)).
% 2. [Step 5‑6] Issue: incorrect cancellation of η‑terms. Re‑derived the
%    inequality correctly, leading to (P6)–(P7) with the factor $(1-\eta)$.
% 3. [Step 7] Issue: unjustified $1-\eta^{2}\ge 1-\beta$. Replaced with a
%    proper use of the step‑size condition $\eta\bar L\le 2(1-\beta)$.
% 4. [Step 8] Issue: unsupported inequality. Replaced by the standard
%    $(a+b)^{2}\ge\frac12 a^{2}-b^{2}$ and used it in (P9).
% 5. [Step 11] Issue: Lemma L3 gave a wrong lower bound. Provided a correct
%    upper bound (Lemma \ref{lem:exp-smooth}) with a full proof.
% 6. Added Lemma \ref{lem:momentum-bound} to control the momentum buffer.
% 7. Introduced a Lyapunov function (P13) to combine the distance and
%    momentum terms, yielding a clean contraction (Step 10).
% 8. Derived the explicit contraction factor $\rho$ and connected it to the
%    step‑size condition, completing the rigorous linear‑rate proof.

\end{document}